\documentclass[12pt,a4paper]{article}

\usepackage[utf8]{inputenc}
\usepackage[T1]{fontenc}
\usepackage[brazil]{babel}
\usepackage{lmodern}
\usepackage[margin=2.5cm]{geometry}
\usepackage{setspace}
\usepackage{graphicx}
\usepackage{booktabs}
\usepackage{tabularx}
\usepackage{xcolor}
\usepackage{tikz}
\usetikzlibrary{positioning,arrows.meta,fit,backgrounds}
\usepackage[hidelinks]{hyperref}

\onehalfspacing
\setlength{\parindent}{1.25cm}

\title{\textbf{Telessaúde segura no interior:}\\
criptografia pós-quântica com satélites LEO e computação de borda}
\author{Cauã Lira\\
\small Licenciatura em Computação -- Universidade Federal Rural de Pernambuco (UFRPE)\\
\small Disciplina: Internet do Futuro -- Profa. Jeisa Pereira de Oliveira Domingues\\
\small 2ª VA -- Etapa 2: Proposta de projeto (teórico)}
\date{Novembro de 2026}

\begin{document}

\maketitle

\begin{abstract}
\noindent
Este documento apresenta uma proposta teórica para proteger os dados de pacientes atendidos por telessaúde em lugares afastados, como o sertão de Pernambuco e comunidades ribeirinhas da Amazônia. A ideia junta três tecnologias vistas na disciplina: a criptografia pós-quântica, que protege os dados contra os computadores quânticos do futuro; os satélites de órbita baixa (LEO), que levam internet até onde não chega fibra óptica; e a computação de borda, que coloca um pequeno servidor dentro do posto de saúde para fazer o trabalho pesado perto do paciente. O texto explica o problema, o papel de cada tecnologia, um esboço da arquitetura em camadas, alguns cenários de uso e uma discussão honesta sobre benefícios, limitações e viabilidade.

\medskip
\noindent\textbf{Palavras-chave:} criptografia pós-quântica; satélites LEO; computação de borda; telessaúde; segurança de dados.
\end{abstract}

% =========================================================
\section{O problema que será investigado}

A telessaúde já é uma realidade no Brasil e, desde 2022, é autorizada por lei em todo o território nacional \cite{brasil2022telessaude}. Ela é especialmente importante para quem mora longe dos grandes centros, porque permite que um paciente em um posto de saúde no interior seja atendido por um especialista que está a centenas de quilômetros de distância. O problema é que muitos desses lugares não têm fibra óptica nem um sinal de celular confiável, e a forma mais prática de levar internet até lá hoje é por satélite.

Junto com a conexão vem uma preocupação que pouca gente enxerga: a segurança dos dados de saúde. Um prontuário, um exame de coração ou um ultrassom de pré-natal são informações que continuam sendo sensíveis pelo resto da vida da pessoa. A Lei Geral de Proteção de Dados (LGPD) inclusive trata esse tipo de informação como dado sensível, que precisa de cuidado redobrado \cite{brasil2018lgpd}. Hoje esses dados viajam protegidos pela criptografia atual, como RSA e curvas elípticas, que funciona bem contra os computadores que temos agora.

Acontece que essa proteção tem prazo de validade. Quando os computadores quânticos ficarem fortes o suficiente, eles vão conseguir quebrar a criptografia atual \cite{bernstein2017pqc}. E o perigo não começa só nesse dia. Já existe o ataque conhecido como \emph{``colha agora, descriptografe depois''}, em que alguém copia os dados criptografados hoje e guarda para abrir no futuro, quando tiver acesso a um computador quântico. Mosca \cite{mosca2018cybersecurity} explica isso de um jeito simples: se o tempo que um dado precisa ficar em segredo, somado ao tempo que levamos para trocar nossa criptografia, for maior do que o tempo até surgir um computador quântico capaz, então já estamos atrasados. Dados de saúde, que precisam ficar em segredo por décadas, se encaixam exatamente nessa situação.

Por isso, o problema investigado nesta proposta é o seguinte: \textbf{como garantir que os dados de pacientes atendidos por telessaúde em regiões remotas continuem protegidos hoje e daqui a vinte ou trinta anos, mesmo usando uma conexão via satélite e equipamentos simples dentro do posto de saúde?}

% =========================================================
\section{Objetivos da proposta}

O objetivo geral é propor uma arquitetura de telessaúde para regiões remotas que proteja os dados dos pacientes contra ameaças atuais e futuras, juntando a criptografia pós-quântica com os satélites de órbita baixa e a computação de borda.

Para chegar nisso, a proposta tem alguns objetivos mais específicos. O primeiro é usar a criptografia pós-quântica para proteger tanto a troca de chaves quanto as assinaturas digitais dos dados de saúde, desde o posto até o hospital de referência. O segundo é usar os satélites LEO para dar uma conexão com velocidade e atraso parecidos com os da internet da cidade, que é o que permite fazer uma teleconsulta em vídeo sem travar. O terceiro é usar a computação de borda para tirar o peso da criptografia dos aparelhos mais simples do posto, além de manter o atendimento funcionando quando o sinal do satélite cair. Por fim, a proposta quer discutir com sinceridade se essa combinação é viável no Brasil de hoje, olhando para custo, energia, dependência de empresas privadas e maturidade das tecnologias.

% =========================================================
\section{As tecnologias e o papel de cada uma}

\subsection{Criptografia pós-quântica (tema 17 -- Cauã)}

A criptografia pós-quântica é um conjunto de novos algoritmos que rodam nos computadores e celulares comuns de hoje, mas que foram pensados para resistir tanto aos computadores normais quanto aos quânticos \cite{bernstein2017pqc}. Em 2024 o NIST, que é o órgão de padronização dos Estados Unidos e uma referência mundial na área, publicou os primeiros padrões oficiais: o ML-KEM, usado para combinar uma chave secreta entre duas partes \cite{nist2024fips203}, e o ML-DSA, usado para assinar digitalmente um documento e provar que ele não foi alterado \cite{nist2024fips204}.

Na nossa arquitetura, ela é a \textbf{camada de proteção}. Ela entra em dois momentos. Na hora de abrir a conexão entre o posto e o hospital, o ML-KEM combina a chave que vai proteger tudo o que passar pelo satélite. E na hora de guardar um exame, o ML-DSA assina o arquivo, para que o médico do outro lado tenha certeza de que aquele exame veio mesmo daquele posto e não foi mexido no caminho. Como esses algoritmos ainda são novos, a proposta usa o chamado \emph{modo híbrido}, que junta o algoritmo clássico com o pós-quântico. Assim, se um dos dois tiver algum defeito descoberto no futuro, o outro continua segurando a proteção.

\subsection{Satélites de órbita baixa -- LEO (tema 8 -- Alexander)}

Os satélites LEO ficam bem mais perto da Terra do que os satélites de internet tradicionais. Enquanto um satélite geoestacionário fica a cerca de 36 mil quilômetros de altura, os de órbita baixa ficam entre algumas centenas e pouco mais de mil quilômetros. Por estarem tão perto, o atraso da conexão cai muito, e por isso eles são lançados em grandes constelações com centenas ou milhares de satélites \cite{delportillo2019leo}. Medições feitas em redes desse tipo mostram atrasos na faixa de dezenas de milissegundos e velocidades suficientes para chamadas de vídeo \cite{michel2022starlink}.

Na nossa arquitetura, os satélites LEO são a \textbf{camada de transporte}. São eles que ligam o posto de saúde isolado ao resto da internet. Um ponto importante é que os dados já saem do posto protegidos pela criptografia pós-quântica. Ou seja, a empresa dona do satélite só transporta os pacotes, mas não consegue ler o que tem dentro deles.

\subsection{Computação de borda -- Edge Computing (tema 2 -- Daniel)}

A computação de borda é a ideia de colocar capacidade de processamento perto de onde os dados nascem, em vez de mandar tudo para uma nuvem distante \cite{shi2016edge}. Satyanarayanan \cite{satyanarayanan2017emergence} mostra que isso ajuda a diminuir o atraso, economizar internet e continuar funcionando mesmo quando a conexão com a nuvem cai.

Na nossa arquitetura, a borda é a \textbf{camada de processamento local}. Ela é um pequeno servidor, do tamanho de um mini computador, instalado dentro do próprio posto de saúde. Esse servidor resolve uma limitação da criptografia pós-quântica que apareceu na Etapa 1: aparelhos simples, como um eletrocardiógrafo portátil ou um tablet antigo, nem sempre têm força para rodar esses algoritmos mais pesados. Com a borda, esses aparelhos só precisam mandar os dados pela rede interna do posto, e o servidor de borda faz o trabalho pesado de criptografar, assinar e enviar. Além disso, ele guarda uma cópia local dos atendimentos, o que permite que o posto continue trabalhando durante uma chuva forte que derrube o sinal do satélite.

\subsection{Resumo dos papéis}

A Tabela~\ref{tab:papeis} resume o que cada tecnologia faz na proposta e qual limitação ela ajuda a resolver.

\begin{table}[h]
\centering
\caption{Papel de cada tecnologia na arquitetura proposta.}
\label{tab:papeis}
\small
\begin{tabularx}{\textwidth}{>{\raggedright\arraybackslash}p{3.2cm} >{\raggedright\arraybackslash}X >{\raggedright\arraybackslash}X}
\toprule
\textbf{Tecnologia} & \textbf{Papel na arquitetura} & \textbf{O que ela resolve} \\
\midrule
Criptografia pós-quântica & Protege a troca de chaves e assina os exames, do posto até o hospital. & Impede que os dados sejam abertos no futuro por um computador quântico. \\
\addlinespace
Satélites LEO & Transporta os dados entre o posto isolado e a internet. & Leva conexão rápida e com pouco atraso a lugares sem fibra. \\
\addlinespace
Computação de borda & Faz a criptografia pesada no lugar dos aparelhos simples e guarda cópia local. & Tira o peso dos aparelhos fracos e mantém o posto funcionando sem sinal. \\
\bottomrule
\end{tabularx}
\end{table}

% =========================================================
\section{Esboço inicial da arquitetura}

A arquitetura foi pensada em quatro camadas, que seguem o caminho do dado desde o paciente até o especialista. A Figura~\ref{fig:arquitetura} mostra esse caminho.

\begin{figure}[h]
\centering
\begin{tikzpicture}[
  font=\small,
  box/.style={draw, rounded corners, align=center, minimum width=3.0cm, minimum height=1.0cm, fill=white},
  layer/.style={draw, dashed, rounded corners, inner sep=6pt},
  arrow/.style={-{Stealth[length=2.5mm]}, thick}
]
% Camada 1
\node[box] (eq) {Aparelhos do posto\\(ECG, tablet, ultrassom)};
% Camada 2
\node[box, right=1.2cm of eq, fill=green!8] (edge) {Servidor de borda\\cripto pós-quântica\\+ cópia local};
% Camada 3
\node[box, right=1.2cm of edge, fill=blue!6] (sat) {Antena + satélites\\LEO + estação\\terrestre};
% Camada 4
\node[box, below=1.6cm of sat, fill=orange!10] (cloud) {Hospital de referência\\/ nuvem de telessaúde};
\node[box, left=1.2cm of cloud] (med) {Médico\\especialista};

\draw[arrow] (eq) -- node[above]{\scriptsize Wi-Fi} node[below]{\scriptsize interno} (edge);
\draw[arrow] (edge) -- node[above]{\scriptsize túnel} node[below]{\scriptsize híbrido} (sat);
\draw[arrow] (sat) -- node[right, align=left]{\scriptsize dados ainda\\\scriptsize cifrados} (cloud);
\draw[arrow] (cloud) -- node[above]{\scriptsize acesso} (med);

\begin{scope}[on background layer]
\node[layer, fit=(eq), label={[font=\scriptsize\bfseries]above:Camada 1: Dispositivos}] {};
\node[layer, fit=(edge), label={[font=\scriptsize\bfseries]above:Camada 2: Borda}] {};
\node[layer, fit=(sat), label={[font=\scriptsize\bfseries]above:Camada 3: Transporte}] {};
\node[layer, fit=(cloud)(med), label={[font=\scriptsize\bfseries]below:Camada 4: Central / Nuvem}] {};
\end{scope}
\end{tikzpicture}
\caption{Esboço da arquitetura proposta, do aparelho no posto até o especialista. Os dados saem da borda já protegidos e só voltam a ser abertos no hospital de referência.}
\label{fig:arquitetura}
\end{figure}

\paragraph{Camada 1 -- Dispositivos.} É onde o atendimento acontece. Aqui ficam os aparelhos usados pela equipe do posto, como um eletrocardiógrafo portátil, um ultrassom de bolso ou o tablet do agente de saúde. Esses aparelhos se ligam ao servidor de borda pela rede Wi-Fi interna do posto, que é protegida com uma senha forte e só aceita aparelhos cadastrados. Eles não precisam rodar criptografia pós-quântica, o que permite aproveitar equipamentos simples ou mais antigos.

\paragraph{Camada 2 -- Borda.} É o coração da proposta. O servidor de borda recebe os dados dos aparelhos, organiza o atendimento e guarda uma cópia local. Antes de mandar qualquer coisa para fora, ele assina cada exame com ML-DSA e abre um túnel seguro com o hospital usando a troca de chaves híbrida, que junta um algoritmo clássico com o ML-KEM. Ele também pode comprimir arquivos grandes, como imagens de ultrassom, para gastar menos internet. Se o sinal cair, ele segura os dados numa fila e envia tudo quando a conexão voltar.

\paragraph{Camada 3 -- Transporte.} É formada pela antena instalada no telhado do posto, pela constelação de satélites LEO e pela estação terrestre que liga os satélites à internet. Essa camada só leva os pacotes de um lado para o outro. Como eles já estão protegidos, nem a empresa do satélite nem alguém que intercepte o sinal no meio do caminho consegue ler o conteúdo.

\paragraph{Camada 4 -- Central / Nuvem.} É o hospital de referência ou o núcleo de telessaúde do estado. Lá, um servidor abre o túnel, confere a assinatura de cada exame e guarda o prontuário já criptografado. O médico especialista acessa o sistema por essa camada para fazer a teleconsulta ou dar o laudo. Como os prontuários ficam guardados por muitos anos, eles também são armazenados com criptografia pós-quântica, e não só durante a viagem.

% =========================================================
\section{Contexto de aplicação e cenários de uso}

O contexto escolhido são as Unidades Básicas de Saúde em áreas rurais e isoladas, como distritos do sertão pernambucano, ilhas e comunidades ribeirinhas da Amazônia. São lugares onde a fibra óptica não chega, onde o sinal de celular é fraco ou inexistente, e onde o especialista mais próximo pode estar a horas de estrada ou de barco. A seguir estão quatro cenários que mostram como a arquitetura funciona no dia a dia.

\paragraph{Cenário 1 -- Teleconsulta com cardiologista.} Seu José, de 67 anos, chega ao posto de um distrito do sertão sentindo uma dor no peito. A técnica de enfermagem faz um eletrocardiograma com um aparelho portátil, que manda o exame para o servidor de borda pelo Wi-Fi do posto. O servidor assina o exame, abre o túnel protegido e envia tudo pelo satélite. Em poucos minutos, um cardiologista no Recife abre o exame, confere a assinatura e conversa por vídeo com a equipe do posto para decidir se Seu José precisa ser levado ao hospital.

\paragraph{Cenário 2 -- Pré-natal na comunidade ribeirinha.} Uma gestante faz um ultrassom com um aparelho de bolso durante a visita mensal da equipe de saúde. O arquivo é grande, então o servidor de borda comprime e envia no fim do dia, quando o posto tem menos uso de internet. O exame fica guardado no prontuário dela no hospital de referência, protegido por muitos anos, já que esse tipo de informação continua sendo íntima e sensível ao longo de toda a vida da mãe e da criança.

\paragraph{Cenário 3 -- Chuva forte e queda do sinal.} Durante uma tempestade, o sinal do satélite cai por algumas horas. O posto não para de atender, porque os aparelhos continuam conversando com o servidor de borda, que guarda tudo numa fila. Quando o sinal volta, o servidor envia os atendimentos acumulados, sempre com a mesma proteção pós-quântica.

\paragraph{Cenário 4 -- Tentativa de ``colha agora, descriptografe depois''.} Imagine que alguém consiga copiar o tráfego que passa pela estação terrestre do satélite e guarde esses dados por anos, esperando um computador quântico. Na arquitetura proposta, essa cópia não serve para nada, porque a chave que protegeu a conexão foi combinada com ML-KEM, que foi feito justamente para resistir a esse tipo de ataque.

% =========================================================
\section{Discussão: benefícios, limitações e viabilidade}

\subsection{Benefícios}

O maior benefício é que a proposta protege hoje um dado que precisa continuar seguro por décadas, em vez de esperar o computador quântico aparecer para só então correr atrás do prejuízo. Outro ponto forte é que as três tecnologias se ajudam de verdade: o satélite leva a conexão, a borda tira o peso da criptografia dos aparelhos simples e mantém o posto funcionando sem sinal, e a criptografia pós-quântica garante que nada do que passa pelo satélite possa ser aberto no futuro. Também vale destacar que os dados não ficam nas mãos da empresa do satélite, já que ela só transporta pacotes que não consegue ler. Por fim, a proposta ajuda a cumprir a LGPD, que exige cuidado especial com dados de saúde \cite{brasil2018lgpd}.

\subsection{Limitações}

A primeira limitação é o tamanho das chaves. Na criptografia de curvas elípticas usada hoje, a chave pública tem 32 bytes. No ML-KEM-768, ela passa para 1.184 bytes, e a assinatura do ML-DSA-65 chega a 3.309 bytes \cite{nist2024fips203,nist2024fips204}. Num link via satélite, onde cada megabyte pode custar caro, isso pesa um pouco mais na abertura de cada conexão. Por outro lado, estudos que testaram esses algoritmos dentro do protocolo TLS mostram que, com uma conexão razoável, o atraso extra é pequeno perto do tempo total da comunicação \cite{paquin2020benchmarking}.

A segunda limitação é a dependência de constelações privadas, que foi exatamente a pergunta do seminário sobre satélites LEO. Se a empresa dona do satélite mudar o preço, cortar o serviço numa região ou sair do país, os postos ficam sem conexão. A criptografia protege o conteúdo, mas não protege contra a falta de serviço.

A terceira limitação é física e prática. O sinal do satélite sofre com chuva forte, o servidor de borda precisa de energia estável, o que nem sempre existe em áreas isoladas, e alguém precisa saber dar manutenção nesses equipamentos. Além disso, a criptografia pós-quântica ainda é nova na prática, e a forma de implementar esses algoritmos sem deixar brechas ainda está sendo amadurecida pela comunidade \cite{bernstein2017pqc}.

\subsection{Viabilidade}

Do ponto de vista técnico, a proposta é viável já hoje. Os algoritmos ML-KEM e ML-DSA estão padronizados desde 2024 \cite{nist2024fips203,nist2024fips204}, e já existem bibliotecas de código aberto e navegadores comuns que usam a troca de chaves híbrida no dia a dia. Os satélites LEO já oferecem internet comercial em áreas rurais do Brasil, com desempenho suficiente para chamadas de vídeo \cite{michel2022starlink}, e o servidor de borda pode ser um mini computador comum, sem nenhum equipamento especial.

O maior desafio está no lado financeiro e humano. Cada posto precisaria de uma antena, de um pequeno servidor, de energia confiável (talvez com placa solar e bateria) e de treinamento para a equipe. Por isso, o caminho mais realista seria começar com um projeto piloto em poucos postos, de preferência em parceria com o núcleo de telessaúde do estado e com uma universidade, medir os resultados e só depois pensar em ampliar. Para diminuir a dependência de uma única empresa, a arquitetura também deixa espaço para trocar o satélite por fibra ou rádio quando essas opções chegarem à região, sem mexer nas camadas de borda e de proteção.

% =========================================================
\section{Considerações finais}

Esta proposta mostrou que a criptografia pós-quântica não precisa ficar restrita a bancos e grandes empresas. Quando ela é combinada com os satélites de órbita baixa e com a computação de borda, dá para levar telessaúde segura até lugares onde a internet mal chega, protegendo os dados dos pacientes não só hoje, mas pelas próximas décadas. As três tecnologias se completam: cada uma cobre um ponto fraco da outra. Ainda existem desafios reais, principalmente de custo, energia e dependência de empresas privadas, mas a base técnica já existe. Como trabalho futuro, a proposta poderia ser fortalecida com o modelo Zero Trust (tema 9), para controlar melhor quem acessa cada prontuário, e com a identidade descentralizada (tema 4), para dar ao paciente mais controle sobre os próprios dados.

\bibliographystyle{abbrv}
\bibliography{refs}

\end{document}
